%% ============================================================================
%% preamble/common.tex -- packages and macros shared by both entry files.
%% Input AFTER \documentclass.  Nothing here may depend on the document class:
%% the same section files are compiled under acmart and under IEEEtran.
%% Requires preamble/switches.tex to have been read (before \documentclass).
%% No spacing commands (\vspace, changed skips, ...) anywhere: two of the four
%% venues reject papers that squeeze the layout.
%% ============================================================================
\makeatletter
\RequirePackage{etoolbox}
\RequirePackage{graphicx}
\RequirePackage{xcolor}
\RequirePackage{booktabs}
\RequirePackage{array}
\RequirePackage{xspace}

\newcommand{\PaperWarn}[1]{\typeout{[paper] WARNING: #1}}

% ---- the PDF records nothing about the machine that built it: no file names of the included figures and
% no engine banner (pdfTeX writes both by default; other engines do not know this parameter)
\ifdefined\pdfsuppressptexinfo \pdfsuppressptexinfo=-1 \fi

% ---- colours (dark enough to survive grey-scale printing) -----------------------
\definecolor{PaperTBD}{RGB}{176,58,0}     % placeholders
\definecolor{PaperStub}{RGB}{64,64,64}    % one-sentence paragraph stubs
\definecolor{PaperTag}{RGB}{110,110,110}  % paragraph ids
\definecolor{PaperFill}{RGB}{140,140,140} % filler text
\definecolor{PaperBox}{RGB}{90,90,90}     % empty figure and table frames

% ============================================================================
% 1. Placeholders
%    \TBD{<id>}{<what it waits for>}   prints  [TBD-<id>: <what it waits for>]
%    \citeTBD{<source>}                prints  [ref: <source>]   (a reference still to be added to refs.bib)
%    One id per open item, the same id wherever the item appears.  Classes:
%      D<n> hardware design items   E<n> experiments   P performance-data source
%      C contribution structure     N names/wording    W text still to be written
%    \PaperTBD = show | blank | none (preamble/switches.tex):
%      show   visible, coloured, listed in the index of open placeholders
%      blank  same layout (as a build without paragraph tags), invisible ink.  The words are
%             still in the PDF text layer: use it to check page counts, never for a submission
%      none   removed, together with the space in front of it
% ============================================================================
\newcommand{\PaperTBDIndexSkip}{W}        % comma list of classes left out of the printed index
\newcommand{\pt@tbdtext}[2]{[TBD-#1: #2]}
\gdef\pt@curpar{--}
\DeclareRobustCommand{\TBD}[2]{%
  \ifcase\PaperTBDmode
    \pt@index{#1}{#2}%
    \begingroup\sffamily\color{PaperTBD}\pt@tbdtext{#1}{#2}\endgroup
  \or
    \begingroup\sffamily\color{white}\pt@tbdtext{#1}{#2}\endgroup
  \else
    \unskip
  \fi}
\DeclareRobustCommand{\citeTBD}[1]{%
  \ifcase\PaperTBDmode
    \pt@index{REF}{#1}%
    \begingroup\sffamily\color{PaperTBD}[ref: #1]\endgroup
  \or
    \begingroup\sffamily\color{white}[ref: #1]\endgroup
  \else
    \unskip
  \fi}
% index of open placeholders: one line per occurrence, written through the .aux file
\newcommand{\pt@index}[2]{%
  \ifPaperTBDIndex
    \addtocontents{tbd}{\protect\PaperTBDEntry{#1}{#2}{\pt@curpar}{\thepage}}%
  \fi}
\newcommand{\PaperTBDEntry}[4]{% id, text, paragraph id, page
  \begingroup
  \def\pt@skip{0}%
  \def\do##1{\ifstrequal{##1}{#1}{\def\pt@skip{1}}{}}%
  \expandafter\docsvlist\expandafter{\PaperTBDIndexSkip}%
  \ifnum\pt@skip=0
    \par\noindent\hangindent=1em\hangafter=1
    {\sffamily\color{PaperTBD}TBD-#1}\quad #3, p.\,#4: #2\par
  \fi
  \endgroup}
\newcommand{\PrintTBDIndex}{%
  \ifPaperTBDIndex
    \clearpage
    \section*{Open placeholders in this draft}%
    {\footnotesize
     \noindent This list is generated from the \texttt{\string\TBD} marks in the text and is not part
     of the paper.  Each line gives the placeholder id, the paragraph id and page where it occurs, and
     what it waits for.  Marks of class \PaperTBDIndexSkip{} (text still to be written) are not listed.\par
     \pt@readtbd}%
  \fi}
% (the kernel's \@starttoc cannot be used: amsart-based classes give it a different signature)
\newcommand{\pt@readtbd}{%
  \begingroup\makeatletter\@input{\jobname.tbd}\endgroup
  \if@filesw
    \newwrite\tf@tbd
    \immediate\openout\tf@tbd=\jobname.tbd\relax
  \fi}

% ============================================================================
% 2. Paragraph units
%    Every paragraph starts with a comment line  %% P:<id>  (read by check_align.py)
%    followed by \ptag{<id>}, which prints the id in small type when \PaperParTags is on.
%    \Intent{...}   one-sentence stub: what the paragraph will say (delete when written)
%    \fillto{<n>}   with \PaperFill=yes, prints <n> filler words: the paragraph's word budget
%    \parhead{...}  run-in paragraph heading that works in both classes
%    \ProtoStatus{...}  the closing "Prototype status:" sentence of a microarchitecture subsection
% ============================================================================
\DeclareRobustCommand{\ptag}[1]{%
  \gdef\pt@curpar{#1}%
  \ifPaperParTags\leavevmode{\scriptsize\sffamily\upshape\mdseries\color{PaperTag}[#1]}\space\fi
  \ignorespaces}
\DeclareRobustCommand{\Intent}[1]{\ifPaperFill\else{\itshape\color{PaperStub}#1}\fi}
\DeclareRobustCommand{\parhead}[1]{\textbf{#1}\space\ignorespaces}
\DeclareRobustCommand{\ProtoStatus}[1]{%   (nothing is printed while the sentence is only a hidden placeholder)
  \ifPaperFill\else\setbox\z@\hbox{#1}\ifdim\wd\z@>\z@ \textit{Prototype status:} #1\fi\fi}

% filler words (their mean length matches the running text of the earlier English draft)
\newcount\pt@fi \newcount\pt@fk \newcount\pt@fK
\global\pt@fK=0
\begingroup
\def\do#1{\expandafter\gdef\csname pt@fw@\the\pt@fK\endcsname{#1}\global\advance\pt@fK by 1 }
\docsvlist{filler,text,stands,in,for,a,paragraph,that,has,not,been,written,yet,and,it,has,exactly,
  the,number,of,words,that,the,page,budget,assigns,to,this,paragraph,so,that,the,page,count,of,this,
  build,approximates,what,the,finished,manuscript,will,require,nothing,in,these,sentences,describes,
  the,architecture,the,evaluation,or,any,result,and,every,word,of,it,disappears,when,the,real,
  paragraph,is,written,the,vocabulary,deliberately,resembles,ordinary,technical,prose,with,a,few,
  longer,expressions,such,as,implementation,characteristics,representation,and,microarchitecture,
  because,line,breaking,and,hyphenation,depend,on,word,length,a,reader,should,ignore,the,wording,
  and,look,only,at,how,much,space,the,surrounding,section,occupies,relative,to,its,figures,tables,
  and,captions,when,the,budget,changes,the,numbers,attached,to,the,individual,paragraphs,change,
  with,it,and,the,pagination,follows,without,any,further,manual,intervention}
\endgroup
\global\pt@fk=0
\newcommand{\pt@filler}[1]{%
  \begingroup\color{PaperFill}\pt@fi=0
  \loop\ifnum\pt@fi<#1\relax
    \csname pt@fw@\the\pt@fk\endcsname\space
    \global\advance\pt@fk by 1
    \ifnum\pt@fk=\pt@fK \global\pt@fk=0 \fi
    \advance\pt@fi by 1
  \repeat
  \unskip.\endgroup}
\DeclareRobustCommand{\fillto}[1]{\ifPaperFill\pt@filler{#1}\fi}

% ============================================================================
% 3. Figure and table slots
%    \FigSlot{<label>}{<what it shows>}{<status>}{<caption>}
%    \TabSlot{<label>}{<what it lists>}{<status>}{<caption>}
%    Width class and body height come from preamble/float-budget.tex, so that an
%    empty slot takes exactly the space the page budget gives it.  Widths are the
%    drawing widths (3.33 in single column, 7.0 in full width) in both classes;
%    a finished figure is included at its natural size and is never scaled, so
%    that lettering drawn at 8 pt stays at 8 pt.
%    When a figure file exists:  \FloatFile{<label>}{figures/<file>}  in float-budget.tex.
%    <what it shows> is also the alternative text of the figure (\Description in acmart): for a
%    finished figure, write it as a description of what the figure actually shows.
%    When a table is typeset, replace its \TabSlot by a table environment that keeps
%    the same \label and the same  %% P:<id>  marker line, names the slot with
%    \SlotName{<label>} and reads the body from tables/<file>.tex; the body starts with \TableSize.
% ============================================================================
\newlength{\PaperColW}  \setlength{\PaperColW}{3.33in}
\newlength{\PaperFullW} \setlength{\PaperFullW}{7.0in}
\newcommand{\FloatBudget}[4]{\csgdef{pf@w@#1}{#2}\csgdef{pf@h@#1}{#3}\csgdef{pf@src@#1}{#4}}
\newcommand{\FloatFile}[2]{\csgdef{pf@file@#1}{#2}}
\newlength{\pf@W}
\newsavebox{\pf@gbox}
\newcommand{\pf@setwidth}[1]{%
  \ifcsdef{pf@w@#1}{}{\PackageError{paper}{No \string\FloatBudget\space entry for `#1'}{Add it to preamble/float-budget.tex.}}%
  \ifcsstring{pf@w@#1}{F}{\setlength{\pf@W}{\PaperFullW}\def\pf@Wtxt{7.0}}{\setlength{\pf@W}{\PaperColW}\def\pf@Wtxt{3.33}}}
\newcommand{\pf@body}[5]{% label, kind word, index class, what, status
  \ifcsdef{pf@file@#1}{%
    \sbox{\pf@gbox}{\includegraphics{\csuse{pf@file@#1}}}%
    \@tempdima=\wd\pf@gbox \advance\@tempdima by -\pf@W
    \ifdim\@tempdima<0pt \@tempdima=-\@tempdima\fi
    \ifdim\@tempdima>1pt
      \PaperWarn{#1: file is \the\wd\pf@gbox\space wide, the slot is \the\pf@W\space wide; draw it at the slot width, it is not scaled}%
    \fi
    \usebox{\pf@gbox}%
  }{%
    \begingroup\setlength{\fboxsep}{3pt}\setlength{\fboxrule}{0.4pt}\color{PaperBox}%
    \fbox{\parbox[c][\dimexpr \csuse{pf@h@#1}in-2\fboxsep-2\fboxrule\relax][c]{\dimexpr\pf@W-2\fboxsep-2\fboxrule\relax}{%
      \centering\footnotesize\sffamily\upshape
      \textbf{#2 slot \detokenize{#1}}\quad(\pf@Wtxt\,in $\times$ \csuse{pf@h@#1}\,in)\par
      Shows: #4.\par
      Status: #5.}}%
    \endgroup
    \ifPaperTBDIndex\addtocontents{tbd}{\protect\PaperTBDEntry{#3}{#5}{\detokenize{#1}}{\thepage}}\fi
  }}
\providecommand{\Description}[2][]{}       % acmart has it (alternative text); IEEEtran does not
\newcommand{\FigSlot}[4]{%
  \pf@setwidth{#1}\gdef\pt@curpar{#1}%
  \ifcsstring{pf@w@#1}{F}%
    {\begin{figure*}[t]\centering\pf@body{#1}{Figure}{FIG}{#2}{#3}\caption{#4}\label{#1}\Description{#2}\end{figure*}}%
    {\begin{figure}[t]\centering\pf@body{#1}{Figure}{FIG}{#2}{#3}\caption{#4}\label{#1}\Description{#2}\end{figure}}}
\newcommand{\TabSlot}[4]{%
  \pf@setwidth{#1}\gdef\pt@curpar{#1}%
  \ifcsstring{pf@w@#1}{F}%
    {\begin{table*}[t]\centering\caption{#4}\label{#1}\pf@body{#1}{Table}{TAB}{#2}{#3}\end{table*}}%
    {\begin{table}[t]\centering\caption{#4}\label{#1}\pf@body{#1}{Table}{TAB}{#2}{#3}\end{table}}}
% \SlotName{<label>}: what \FigSlot and \TabSlot do themselves -- the index of open placeholders names the
% float in which a placeholder stands.  Written once at the top of a float that is typeset by hand.
\newcommand{\SlotName}[1]{\gdef\pt@curpar{#1}}
% \TableSize: the type size inside typeset tables, 8 pt in every page format (\footnotesize is 8 pt where
% the body is set in 10 pt and 7 pt where it is set in 9 pt; there \small is 8 pt).  The body size is
% read once at \begin{document}: inside a float the current size is not a guide (IEEEtran switches to
% \footnotesize in every table float).
\newcommand{\TableSize}{\footnotesize}
\AtBeginDocument{\begingroup\normalsize\ifdim\f@size\p@<9.5\p@ \gdef\TableSize{\small}\fi\endgroup}
% geometry of the text block, read by page_report.py (written at the end: acmart becomes two-column in \maketitle)
\AtEndDocument{%
  \@tempdima=1in \advance\@tempdima by \voffset \advance\@tempdima by \topmargin
  \advance\@tempdima by \headheight \advance\@tempdima by \headsep
  \typeout{[paper] geometry: top=\the\@tempdima\space textheight=\the\textheight\space textwidth=\the\textwidth\space columnwidth=\the\columnwidth}}
\AtEndDocument{%
  \ifdim\PaperColW>\columnwidth
    \PaperWarn{single-column slot (\the\PaperColW) is wider than the column (\the\columnwidth)}\fi
  \ifdim\PaperFullW>\textwidth
    \PaperWarn{full-width slot (\the\PaperFullW) is wider than the text block (\the\textwidth)}\fi}

% ============================================================================
% 4. Cross references (the words differ by class; the entry files may redefine them)
% ============================================================================
\newcommand{\figrefname}{Figure}
\newcommand{\tabrefname}{Table}
\newcommand{\secrefname}{Section}
\newcommand{\figref}[1]{\figrefname~\ref{#1}}
\newcommand{\tabref}[1]{\tabrefname~\ref{#1}}
\newcommand{\secref}[1]{\secrefname~\ref{#1}}

\makeatother
%% ============================================================================
%% preamble/meta.tex -- title, names and links: the things that change between a
%% share build and a submission, kept in one place.
%% ============================================================================

% Working title.  A build that prints names puts the project name in front of it; an anonymous
% build does not (double-blind venues ask for a system name and a title that differ from any
% public version: set \PaperTitleCore differently for a submission if the venue asks for it).
\newcommand{\PaperTitleCore}{Keeping GEMM Results Block-Scaled across the Tensor--Vector Hand-off: When Does It Pay Off?}
\newcommand{\PaperTitleCorePDF}{Keeping GEMM Results Block-Scaled across the Tensor-Vector Hand-off: When Does It Pay Off?}
\newcommand{\PaperSysNamePublic}{ScaleNative}  % the project name
\newcommand{\PaperSysNameAnon}{[System]}      % the name used in an anonymous submission
\ifPaperAnon
  \newcommand{\PaperTitle}{\PaperTitleCore}
  \newcommand{\PaperTitlePDF}{\PaperTitleCorePDF}
\else
  % The signed build prints the working title without the project name: "ScaleNative: ...: When Does
  % It Pay Off?" reads badly with two colons and wraps to three lines.  The name is used in the text
  % (\SysName).  To put it in the title, replace the two lines below by
  %   \newcommand{\PaperTitle}{\PaperSysNamePublic: \PaperTitleCore} (and likewise \PaperTitlePDF).
  \newcommand{\PaperTitle}{\PaperTitleCore}
  \newcommand{\PaperTitlePDF}{\PaperTitleCorePDF}
\fi
\newcommand{\PaperShortTitle}{Keeping GEMM Results Block-Scaled}
\newcommand{\PaperVersion}{first full text, 2026-10-06}
\newcommand{\PaperKeywords}{}                 % left empty until the text exists

% System name in the text.  A submission to a double-blind venue must not use the public name.
\newcommand{\SysName}{\ifPaperAnon\PaperSysNameAnon\else\PaperSysNamePublic\fi\xspace}

% Repository.  The public repository (printed by a build that shows names; never by an anonymous
% build).  While this is empty \RepoLink prints a placeholder.
\providecommand{\PaperRepoURL}{https://github.com/fifakaka35357/ScaleNative}
% \RepoLink prints the public URL, an anonymized URL (\PaperRepoAnonURL) or a neutral phrase.
\newcommand{\RepoLink}{%
  \ifcase\PaperRepoMode
    \ifdefempty{\PaperRepoURL}{\TBD{N}{repository URL}}{\expandafter\url\expandafter{\PaperRepoURL}}%
  \or
    \expandafter\url\expandafter{\PaperRepoAnonURL}%
  \else
    [artifact link withheld for double-blind review]%
  \fi}

% Note attached to the title of a share build (never printed in a submission).
\newcommand{\PaperNoticeText}{Working draft (\PaperVersion) under a working title; not peer reviewed, not submitted to any venue.
  Colored brackets mark open items; gray tags name paragraphs; empty frames stand for figures to come.%
  \ifdefempty{\PaperRepoURL}{\ifnum\PaperTBDmode<2 { Sources: \RepoLink}\fi}{ Sources: \RepoLink}}
\endinput

%% ============================================================================
%% preamble/terms.tex -- terminology macros.
%% Use these for every term whose wording must stay the same across the paper or
%% may still change.  Each macro ends in \xspace: write "\TensorUnit produces".
%% The Chinese word used in paper.zh-CN.md is given in the comment.
%% Entries marked UNDECIDED are working choices that the first author has not confirmed.
%% Names that are lettered on the figures (figures/) are written here as they are written there:
%% Tensor unit, Vector unit, Vector core, L2, capsule; the block ``Fast lane'' is ``fast lane'' in running text.
%% ============================================================================

% ---- units of the PE ------------------------------------------------------------
\newcommand{\TensorUnit}{Tensor unit\xspace}          % Tensor 单元
\newcommand{\VectorUnit}{Vector unit\xspace}          % Vector 单元（两个 Vector 核的总称）
\newcommand{\VectorCore}{Vector core\xspace}          % Vector 核（每个 PE 两个）
\newcommand{\Lii}{L2\xspace}                          % L2（PE 内的本地存储）

% ---- the five hand-off placements (B/S = block/scalar; D/M/E = direct/through memory/epilogue)
\newcommand{\armBD}{\mbox{B-D}\xspace}                % 块直通
\newcommand{\armSD}{\mbox{S-D}\xspace}                % 标量直通
\newcommand{\armSE}{\mbox{S-E}\xspace}                % 生产端融合
\newcommand{\armBM}{\mbox{B-M}\xspace}                % 块经存储
\newcommand{\armSM}{\mbox{S-M}\xspace}                % 标量经存储

% ---- recurring terms ------------------------------------------------------------
\newcommand{\blockscaled}{block-scaled\xspace}        % 块缩放
\newcommand{\handoff}{hand-off\xspace}                % 交接
\newcommand{\failvisible}{fail-visible\xspace}        % 失败可见
\newcommand{\capsule}{capsule\xspace}                 % capsule（交接时的数据加元数据）
\newcommand{\sidecar}{scale sidecar\xspace}           % scale sidecar（与 payload 分开存放的块指数）
\newcommand{\FastPath}{fast lane\xspace}              % 快通道（图上的块名是 Fast lane；宏名沿用旧名）

% ---- forms of evidence (a physical result is never called "measured") --------------
\newcommand{\evRTL}{RTL simulation\xspace}            % RTL 仿真
\newcommand{\evModel}{cycle-level model\xspace}       % 周期模型（方法流程图上是留空的槽位 Model results：用哪套模型还没有选）
\newcommand{\evAnalytic}{analytical derivation\xspace} % 解析推导
\newcommand{\evSlice}{synthesis slice\xspace}         % 综合切片
\newcommand{\evPostRoute}{post-route EDA estimate\xspace} % 布线后 EDA 估计

% ---- configuration names (UNDECIDED: neutral English names for the three configurations) --
\newcommand{\cfgAsBuilt}{\textsc{Prototype}\xspace}   % 今天的 RTL
\newcommand{\cfgCorrected}{\textsc{Corrected}\xspace} % 修正后的 PE
\newcommand{\cfgProjected}{\textsc{Projected}\xspace} % 投影的系统

% "heterogeneous" is reserved for a mesh that mixes two PE types; a PE with a Tensor unit
% and Vector cores is not called heterogeneous.
\endinput

%% preamble/float-budget.tex -- width class and body height of every figure and table slot.
%% Source of the numbers: page-budget plan (14 figures + 7 tables); S = single column (3.33 in),
%% F = full width (7.0 in); heights are body heights in inches, caption excluded.
%% \FloatBudget{<label>}{<S|F>}{<height in inches>}{<free note, not printed>}
%% The seven tables are typeset (bodies in tables/); their entries stay as the record of the size each was given.
%% To use a finished figure instead of the empty box:  \FloatFile{<label>}{figures/<file>}

\FloatBudget{fig:handoff}{S}{1.9}{}% Fig. 1
\FloatBudget{fig:workload}{S}{2.1}{}% Fig. 2
\FloatBudget{tab:workload-gen}{S}{1.8}{}% Tab. 1 (typeset: tables/tab-workload-gen.tex; this entry is no longer read)
\FloatBudget{tab:placements}{F}{2.2}{}% Tab. 2 (typeset: tables/tab-placements.tex; no longer read)
\FloatBudget{tab:attribution}{S}{1.3}{}% Tab. 3 (typeset: tables/tab-attribution.tex; no longer read)
\FloatBudget{fig:retain}{S}{2.6}{}% Fig. 3
\FloatBudget{fig:system}{S}{2.1}{}% Fig. 4
\FloatBudget{fig:pe}{F}{2.4}{}% Fig. 5
\FloatBudget{fig:tensor}{F}{1.9}{}% Fig. 6
\FloatBudget{fig:vector}{F}{2.4}{}% Fig. 7
\FloatBudget{fig:flow}{S}{1.5}{}% Fig. 8
\FloatBudget{tab:configs}{S}{1.2}{}% Tab. 4 (typeset: tables/tab-configs.tex; no longer read)
\FloatBudget{tab:mac-share}{S}{1.3}{}% Tab. 5 (typeset: tables/tab-mac-share.tex; no longer read)
\FloatBudget{fig:payoff}{F}{2.3}{}% Fig. 9
\FloatBudget{fig:correct}{S}{2.0}{}% Fig. 10
\FloatBudget{fig:metadata}{S}{2.0}{}% Fig. 11
\FloatBudget{fig:failvis-cost}{S}{1.9}{}% Fig. 12
\FloatBudget{fig:layout}{S}{3.33}{}% Fig. 13 (the die is square and the figure has no margin, so it is as high as it is wide)
\FloatBudget{tab:layout-runs}{S}{1.6}{}% Tab. 6 (typeset: tables/tab-layout-runs.tex; no longer read)
\FloatBudget{fig:system-proj}{S}{2.0}{}% Fig. 14
\FloatBudget{tab:prior}{S}{1.6}{}% Tab. 7 (typeset: tables/tab-prior.tex; no longer read)

%% Finished figures (sources and build script in figures/).  Each PDF is drawn at the width and height
%% budgeted above and is included at its natural size.
\FloatFile{fig:system}{figures/fig-system}% Fig. 4,  3.33 in x 2.1 in
\FloatFile{fig:pe}{figures/fig-pe}% Fig. 5,  7.0 in x 2.4 in
\FloatFile{fig:tensor}{figures/fig-tensor}% Fig. 6,  7.0 in x 1.9 in
\FloatFile{fig:vector}{figures/fig-vector}% Fig. 7,  7.0 in x 2.4 in
\FloatFile{fig:flow}{figures/fig-flow}% Fig. 8,  3.33 in x 1.5 in
\FloatFile{fig:layout}{figures/fig-layout}% Fig. 13, 3.33 in x 3.33 in

\endinput
